% sections/graphs.tex -- Methods B: graphs and integrands (generation, construction, consistency checks)
% ---------------------------------------------------------------------------------------------------------------
\subsection{Graphs and integrands}
\label{sec:methods-graphs}

The pipeline from the Lagrangian to the sampled integrands has three stages, each with its own checks: for each order, the generation
of the graphs, the construction of the subtracted integrand of every graph, and the consistency checks that
tie the two to exact results. All three are automatic and were run here to ten loops.

\subsubsection{Generating the graphs}
\label{sec:methods-generation}

The lists of 1PI graphs of $\Gamma^{(2)}$ and $\Gamma^{(4)}$ are produced by a generator that derives from the one
written for the seven-loop computation of Ref.~\cite{Ribeca2003thesis}, with three improvements: it grows the graphs
by vertex insertion instead of enumerating labelled graphs and screening them by a canonical form, it canonicalizes
them by colour refinement with backtracking, and it takes any list of vertex degrees, the external legs being implicit
in the valences. The class it generates is every connected graph with the given valences, self-loops and bridges
allowed, and the construction is complete: removing a leaf of a spanning tree and re-pairing the half-lines it leaves
is the inverse of an insertion, so that every graph with $V$ vertices descends from one with $V-1$. The class
is grown order by order, each candidate canonicalized and kept in a hash set, and the 1PI graphs are written out at
the end of every order, together with their Nickel indices and weights. The cost is a few hundred insertions per
graph grown, so that it follows the size of the class, growing by a factor of ten to twelve per loop on one core.
Table~\ref{tab:generation} gives the sizes and the times, measured on one core of a cluster node; the
counts through eight loops are those of Ref.~\cite{Ribeca2003thesis} (Table~3.1), the others are new~\check{G1}. Every
complete list satisfies its zero-dimensional sum rule (Sec.~\ref{sec:methods-checks}) through twelve loops for
$\Gamma^{(2)}$ and eleven for $\Gamma^{(4)}$.

Every graph is converted to its Nickel index~\cite{NMB1977} in canonical form---the
lexicographic minimum over vertex labellings---which identifies it across lists and keys its random numbers
(Sec.~\ref{sec:methods-rng}); one-vertex-reducible graphs and graphs with both legs on one vertex are flagged, since
their integrands vanish (Sec.~\ref{sec:methods-construction}), and are not built for $B_L$ and $H_L$.

\begin{table}
\caption{\label{tab:generation}%
Generation of the graph lists: numbers of 1PI graphs without self-loops (one-vertex-reducible ones included), and the
CPU time of the generator on one core of a cluster node (an Intel Xeon Platinum 9242), one list at a time on an
otherwise idle node, the Nickel indices, weights and sum rules included. The eleven- and twelve-loop lists were made
on 96 threads, in 23~s and 275~s of wall time for $\Gamma^{(2)}$ and 1005~s for $\Gamma^{(4)}$ at eleven loops, the
last with 38~GB of memory; their times in the table are core times summed over the threads. Every regenerated list
equals the stored one. The counts of graphs whose integrands survive the pruning (Sec.~\ref{sec:methods-construction})
are in Table~\ref{tab:graphcounts}.}
\begin{ruledtabular}
\setlength{\tabcolsep}{3pt}\footnotesize
\begin{tabular}{lrrrr}
$L$ & $\Gamma^{(4)}$ & time & $\Gamma^{(2)}$ & time \\
\hline
7  & 3986       & 2.1 s & 317    & 0.1 s \\
8  & 26\,540    & 22 s & 1622   & 1.1 s \\
9  & 197\,022   & 4.3 min & 9184   & 9.8 s \\
10 & 1\,593\,254 & 53 min & 59\,614 & 98 s \\
11 & 13\,975\,511 & 14.1 h & 428\,164 & 22 min \\
12 & --- & --- & 3\,393\,754 & 4.2 h \\
\end{tabular}
\end{ruledtabular}
\end{table}

\paragraph{Weights.}
Each graph carries the weight $S_G=T_G(N)\,E!/|\mathrm{Aut}\,G|$, where $E$ is its number of legs, $\mathrm{Aut}\,G$ its
automorphism group at the level of half-edges with unlabelled legs, and $T_G(N)$ its $O(N)$ factor, obtained by
contracting the vertex tensors $(\delta_{ij}\delta_{kl}+\delta_{ik}\delta_{jl}+\delta_{il}\delta_{jk})/3$ over the $3^V$
pairings at the $V$ vertices and projecting on the external tensor, so that $T_G(1)=1$. The factor $E!/|\mathrm{Aut}\,G|$
includes the sum over the labellings of the legs, which is legitimate at zero external
momenta~\cite{Ribeca2003thesis}. The integrals do not depend on $N$: since every graph carries its polynomial
$T_G(N)$, one run gives the coefficients for all $N$.

\subsubsection{Constructing the integrands}
\label{sec:methods-construction}

Every polynomial in Eq.~\eqref{eq:forestterm} is the Kirchhoff polynomial of a minor of $G$, possibly with two vertices
identified, and factorizes into the Kirchhoff polynomials of its blocks, which we call atoms. Expanding all forest terms
into products of half-integer powers of atoms, with binomial coefficients that are dyadic rationals and hence exact in
binary floating point, and adding the coefficients of identical products, removes every combination of forest terms
that cancels identically; we call this step the pruning. Each graph becomes
\begin{equation}
f_G(x)=\sum_k c_k\,x_{\mathrm{ins},k}\prod_a P_a(x)^{p_{ka}/2},
\label{eq:integrand}
\end{equation}
with atoms $P_a$, integer $p_{ka}$ and exact $c_k$; terms with the same atoms, which differ only by the insertions, are
grouped. The result of the expansion is a small program: the list of atoms (each a minor of $G$, evaluated at run time
by a spanning-tree pass and a factorization of its loop matrix), the list of terms (coefficient, insertion monomial,
atom exponents), and the tables of the sampling density of Sec.~\ref{sec:methods-tropical}. One generic kernel
evaluates any such program (Sec.~\ref{sec:methods-kernel}).

The pruning removes up to two thirds of the raw terms. At six loops a four-point graph has on average 250 forests,
which leave 116 terms in 22 groups with 19 distinct atoms (at most 55); at seven loops, in a sample of 300 four-point
graphs, 583 forests, 219 terms, 39 groups and 26 atoms (at most 6128, 766, 126 and 85). The largest programs are
those of two-point graphs with many nested two-point parts: at nine loops the chain of eight bubbles of $C_9$ has
$3.4\times10^5$ forests and 2328 atoms, against at most 269 atoms for $H_8$ and 203 for $B_8$. The expansion is done
once per graph on the host, in parallel over the graphs; its cost grows faster than the number of graphs
(Table~\ref{tab:graphcounts}: about \buildpc{b8} core-seconds per graph for $B_8$, \buildpc{h9} for $H_9$,
\buildpc{c9} for $C_9$, and from \buildpc{b10} for $B_{10}$ to \buildpc{c10} for $C_{10}$ at ten loops), so that at ten
loops it dominates a job. The built integrands of a list are therefore
written once to a cache file, by a CPU-only job, and loaded by every subsequent run; a run with and without the cache
gives bitwise the same results~\check{I4}.

Graphs that vanish identically come out of the pruning with no terms at all, without any structural rule: graphs with a
cut vertex (one-vertex-reducible graphs), which factorize at zero momentum~\cite{Ribeca2003thesis} into pieces of which
one is a renormalization part at exactly zero momentum, cancelled by its own subtraction; graphs with self-loops,
for $L\ge2$; and, for $H_L$, generalized tadpoles, graphs with both external legs on one vertex, for which $W_G=0$. At
three to seven loops these are 21--30\% of the four-point and 15--17\% of the two-point lists (more for $H_L$); at
eight loops 18\% and 14\%, and at nine loops 12\% of the two-point list. In production runs the one-vertex-reducible
graphs of $B_L$ and $H_L$ are dropped before the expansion; their vanishing, and that of the graphs with self-loops, was
verified explicitly up to five loops~\check{I2}. Those of $C_L$ are expanded in every run, and the pruning leaves no
term in any of them, at every order (independently verified up to seven loops~\check{I2}).

\paragraph{Twins.}
At zero momenta and equal masses, Eq.~\eqref{eq:param} depends on $G$ only through its cycle matroid, so that
2-isomorphic graphs in Whitney's sense~\cite{Whitney1933}, which we call twins, have equal integrals in any
representation~\cite{Ribeca2003thesis}: the 3986 seven-loop four-point graphs have 3018 distinct unsubtracted integrals,
and the 3150 non-vanishing ones fall into 2391 classes. We checked that the subtracted integrands of twins, summed over
the positions of the insertion, agree pointwise under the matroid map of the lines (all 3 pairs at four loops and all 25
non-vanishing pairs at five loops for $B_L$)~\check{I3}. Twins have different $O(N)$ factors, so that a merged class carries the
sum of their weights. Twins were not merged in the runs of Sec.~\ref{sec:results}: every class is one graph.

\subsubsection{Consistency checks}
\label{sec:methods-checks}

The graphs, their weights and their integrands are checked, before any sampling, against results that are exact.
\begin{enumerate}
\item \emph{Zero-dimensional sum rules.} In zero dimensions every integral equals one, so that $\sum_{G}S_G$ over a
complete list is the coefficient of $g^L$ in the corresponding vertex function of the zero-dimensional $O(N)$ model,
known exactly from its moments (Appendix~\ref{app:sumrules}). Summed over all graphs, one-vertex-reducible ones
included, this tests the completeness of a list and its weights without any integral; for lists without self-loops the
bare mass is shifted so that every self-loop cancels. Every list passes it exactly, with the full $O(N)$
polynomial, through ten loops (Table~\ref{tab:sumrules})~\check{G2}.
\item \emph{Independent lists.} Converted to Nickel indices, our lists equal those of Ref.~\cite{Phi4tools2024}
(generated with \textsc{feyngen}~\cite{Borinsky2014feyngen}) as sets, for $\Gamma^{(4)}$ at $L\le7$ and $\Gamma^{(2)}$ at
$L\le8$~\check{G3}, and those of a brute-force generator in our code where that is practical ($\Gamma^{(4)}$ to five loops,
$\Gamma^{(2)}$ to six). Two independent canonicalizers (one in the generator's string format, one on Nickel indices)
agree on every graph, and all graphs of every list are distinct~\check{G1}.
\item \emph{Weights.} The weights and $O(N)$ polynomials equal those of the independent implementation of
Refs.~\cite{Ribeca2003thesis,GuidaRibeca2006} on all 260 graphs with $L\le5$ (its 251 and the nine graphs with self-loops its lists lacked) and on all 47
graphs of its four-loop tables~\check{G4}.
\item \emph{Integrands.} The expansion~\eqref{eq:integrand} agrees pointwise with an independent evaluation of the
forest formula in 100-digit arithmetic (Sec.~\ref{sec:methods-validation})~\check{I1}; vanishing graphs
vanish~\check{I2}, twins agree pointwise~\check{I3}, and the exact two-loop sums $B_2=1/2$, $H_2=1/81$ and $C_2=1/12$
are reproduced to a few $10^{-7}$ by the sampled integrals (Sec.~\ref{sec:results-validation})~\check{N3}.
\item \emph{Large order.} Each new coefficient is compared with its leading large-order estimate
(Sec.~\ref{sec:methods-validation}), a test of the complete chain that catches an error of a few per cent (Sec.~\ref{sec:methods-validation})~\check{S3}.
\end{enumerate}
